\section{引导阶段区域注册}
\label{sec:vma-boot}

在 \texttt{kmain.c} 中，VMA 初始化后立即注册两个引导时建立的内核区域：

\codefile{src/kernel/core/kmain.c}
\begin{lstlisting}[style=cstyle]
vma_init();

if (vma_is_ready()) {
    /* 1. 直接映射区: 内核代码/数据 + 直接映射的物理内存 */
    vma_add(KERNEL_VIRT_OFFSET, vmm_direct_map_end(),
            VMA_READ | VMA_WRITE | VMA_EXEC, "direct-map");

    /* 2. 内核堆区 */
    vma_add(KHEAP_START,
            KHEAP_START + KCONFIG_HEAP_INITIAL_PAGES * VMM_PAGE_SIZE,
            VMA_READ | VMA_WRITE, "kheap");
}
\end{lstlisting}

\begin{itemize}
  \item \texttt{direct-map}：覆盖 \hex{C0000000} 到 \texttt{vmm\_direct\_map\_end()}，
        包含内核代码（需要 EXEC）、只读数据和读写数据。权限给足 \texttt{rwx}，
        更细粒度的分段保护留给将来的 \texttt{.text}/\texttt{.rodata} 拆分。

  \item \texttt{kheap}：覆盖 \hex{E0000000} 开始的初始堆区域。
        堆只做数据 \texttt{rw}，无需 \texttt{x}。
        堆增长时应同步调用 \texttt{vma\_add} 或更新 end 来跟踪新区域。
\end{itemize}

\begin{designbox}
注册 VMA 的时机是"谁建立映射，谁负责注册"：
\begin{itemize}
  \item \texttt{vmm\_init()} 建立直接映射 → \texttt{kmain} 注册 \texttt{direct-map}
  \item \texttt{kmalloc\_init()} 建立堆映射 → \texttt{kmain} 注册 \texttt{kheap}
  \item \texttt{vmm\_alloc\_pages()} 将来可在内部自动调用 \texttt{vma\_add}
\end{itemize}
这比集中在一处注册更不容易遗漏——映射与注册在同一上下文中完成。
\end{designbox}

% ============================================================================

